\documentclass[10pt,a4paper]{amsart}
\usepackage[margin=19mm]{geometry}
\usepackage{amsmath,amssymb,mathtools}
\usepackage[scr=rsfs,cal=euler]{mathalfa}
\usepackage{xfrac}
\usepackage[unicode,hidelinks,bookmarks=false]{hyperref}
\newcommand{\ie}{i.e.\ }
\newcommand{\cf}{cf.\ }
\newcommand{\resp}{respectively\ }
\newcommand{\Subsection}{\subsection}
\providecommand{\nobreakdash}{\nobreak\hbox{-}}
\setlength{\emergencystretch}{2em}
\multlinegap0pt
\usepackage{bm}
\usepackage{bbm}
\usepackage{dsfont}
\usepackage{xspace}
\usepackage{cancel}
\usepackage{mathtools}
\usepackage{amsthm}
\makeatletter
\@ifundefined{defi}{\newtheorem{defi}{Definition}}{}
\@ifundefined{proposition}{\newtheorem{proposition}{Proposition}}{}
\makeatother
\def \crit {2^\star}
\def \eps {\epsilon}
\def \mai {\mu_{\alpha,i}}
\def \nap {\nu_{\alpha,p}}
\def \naq {\nu_{\alpha,q}}
\def \rn {\mathbb{R}^n}
\def \tuai {\tilde{u}_{\alpha,i}}
\def \tvap {\tilde{v}_{\alpha,p}}
\def \ua {u_\alpha}
\def \xai {x_{\alpha,i}}
\def \xaj {x_{\alpha,j}}
\def \yap {y_{\alpha,p}}
\def \yaq {y_{\alpha,q}}
\title[Sharp multiscale control for high order nonlinear equations]{Sharp multiscale control for high order nonlinear equations}
\author{Fr\'ed\'eric Robert}
\date{Source: arXiv:2509.07517v1 (CC BY 4.0)}
\begin{document}
\maketitle

\noindent\emph{Source and licence.} Sharp multiscale control for high order nonlinear equations. Version arXiv:2509.07517v1 (CC BY 4.0), distributed under the Creative Commons Attribution 4.0 International licence, as recorded in the arXiv licence metadata for this exact version and shown on the abstract page https://arxiv.org/abs/2509.07517v1. Full rights notice: licenses/arxiv-2509.07517-CC-BY-4.0.txt.\par\smallskip

\noindent\textit{Editorial scope.} Counted unit: the Proposition whose source label is \texttt{prop:4.4} in the article (source0.tex lines 376--377, source label \texttt{prop:4.4}), delivered together with the article's own complete original proof of it (source0.tex lines 378--424, one \texttt{proof} environment of 47 lines). No mathematics is written, repaired or re-derived: statement and proof are the article's own text line for line. Required context retained uncounted: lem:lip (lines 153--155); hyp:bnd:nrj:pf (lines 175--177); min:nrj:U (lines 215--217); cv:tuai (lines 242--258); lim:d:infty (lines 254--256); prop:3.2 (lines 293--298); ineq:dist:xai:yap (lines 350--371); cv:tvap (lines 365--367); def:Sp (lines 368--370). Every definition, display and label of those lines is retained unchanged. Omitted from this package and not redistributed: 1--152, 156--174, 178--214, 218--241, 257--292, 299--349, 371--375, 425--1973. Nothing inside a retained excerpt is omitted or altered: every retained body is delivered exactly as the article's lines are written, so no omission receipt of any kind accompanies this package. Citations: the retained excerpts carry no citation command at all, so no bibliography entry of the article is transcribed and no reference is redistributed. Reference adaptation: none was needed and none was made; every reference inside the delivered unit resolves inside it, so no unresolved reference or citation is produced. Printed slips kept literal: no printed word, symbol, hypothesis or number of the counted statement or of its proof was altered. Layout and class: the article's own class is \texttt{amsart} with packages including amsmath, amssymb, stmaryrd, amsrefs, amsfonts, hyperref, cleveref, upref; the wrapper is the lane's pinned \texttt{amsart} editorial wrapper at 10pt on A4, which restates the theorem-like environments the retained text needs and adds the article's own macro definitions that the retained text needs (\texttt{\textbackslash crit}, \texttt{\textbackslash eps}, \texttt{\textbackslash mai}, \texttt{\textbackslash nap}, \texttt{\textbackslash naq}, \texttt{\textbackslash rn}, \texttt{\textbackslash tuai}, \texttt{\textbackslash tvap}, \texttt{\textbackslash ua}, \texttt{\textbackslash xai}, \texttt{\textbackslash xaj}, \texttt{\textbackslash yap}, \texttt{\textbackslash yaq}), with extra wrapper packages (bm, bbm, dsfont, xspace, cancel, mathtools). The delivered numbering is the unit's own. Assets: no retained span contains a figure environment, a graphics-inclusion command, a PDF-page inclusion, a table or a TikZ picture; no image file is copied into this package, the package declares no asset dependency and no image file is redistributed with it. Licence: version arXiv:2509.07517v1 of this article is distributed under the Creative Commons Attribution 4.0 International licence, as recorded in the arXiv licence metadata for this exact version and shown on the abstract page https://arxiv.org/abs/2509.07517v1. A full rights notice is delivered at licenses/arxiv-2509.07517-CC-BY-4.0.txt. Characters: every retained excerpt is pure ASCII, so the delivered engine, XeLaTeX with its default Latin Modern text font, reports no missing character.
\begin{equation}\label{lem:lip}
\frac{1}{2}|X-Y|\leq d_g(\hbox{exp}_p(X),\hbox{exp}_p(Y))\leq 2|X-Y|
\end{equation}

and \begin{equation}\label{hyp:bnd:nrj:pf}
\Vert\ua\Vert_{\crit} \leq \Lambda \hbox{ for all }\alpha>0.
\end{equation} 

\begin{equation}\label{min:nrj:U}
\int_{\rn}|V|^{\crit}\, dv_\xi\geq \frac{1}{K(n,k)^{\frac{n}{2k}}}.
\end{equation}

\begin{defi} Given $K\geq 1$, we say that $(H_{K})$ holds if for all $i=1,..,K$ there exists $(\xai)_\alpha\in M$ and there exists $U_i\in D_k^2(\rn)\cap C^{2k}(\rn)-\{0\}$ such that:
\begin{itemize}
\item $\mai:=|\ua(\xai)|^{\frac{-2}{n-2k}}\to 0$ as  $\alpha\to +\infty$;
\item For all $x\in\rn$,
\begin{equation}\label{cv:tuai}
\tuai :=\mai^{\frac{n-2k}{2}}\ua(\hbox{exp}_{\xai}(\mai \cdot))\to U_i\hbox{ in }C^{2k}_{loc}(\rn)
\end{equation}
\item We have that $\Delta_\xi^k U_i=|U_i|^{\crit-2}U_i$ strongly in $\rn$ and weakly in $D_k^2(\rn)$ and
\begin{equation*}
\int_{\rn}|U_i|^{\crit}\, dv_\xi\geq \frac{1}{K(n,k)^{\frac{n}{2k}}}.
\end{equation*}
\item For all $i,j\in \{1,...,K\}$, $i\neq j$, we have that
\begin{equation}\label{lim:d:infty}
\lim_{\alpha\to +\infty}\frac{d(\xai,\xaj)}{\mai}=+\infty.
\end{equation}
\end{itemize}
\end{defi}

\begin{proposition}\label{prop:3.2} There exists $K\geq 1$ such that $(H_K)$ holds and 
\begin{equation}\label{hyp:weak:2}
 R_{\alpha, K}(x)^{\frac{n-2k}{2}}|\ua(x)|\leq C\hbox{ for all }x\in M\hbox{ and }\alpha>0
\end{equation}
 where $R_{\alpha, K}(x):=\min_{i=1,...,K}d_g(x,\xai)$ for all $x\in M$ and $\alpha>0$.
\end{proposition}

\begin{defi} We say that $(\tilde{H}_{0})$ holds if $(H_K)$ holds as in Proposition \ref{prop:3.2}. Given $L\geq 1$, we say that $(\tilde{H}_{L})$ holds if for all $p=1,...,L$ there exists $(\yap)_\alpha\in M$ such that $\nap:=|\ua(\yap)|^{\frac{-2}{n-2k}}\to 0$ as  $\alpha\to +\infty$, and, denoting
$$I_p:=\{i\in\{1,...,K\}\hbox{ such that }d_g(\xai,\yap)=O(\nap)\},$$
there exists $\eps_0>0$ such that for all $i\in\{1,...,K\}$ and $p\in\{1,...,L\}$, we have that
\begin{equation}\label{ineq:dist:xai:yap}
\lim_{\alpha\to+\infty}\frac{d_g(\yap,\xai)}{\mai}=+\infty\hbox{ and }\frac{d_g(\yap,\xai)}{\nap}\geq\eps_0.
\end{equation}
Moreover, for $p,q\in \{1,...,L\}$ such that $p\neq q$, then
either 
\begin{equation}\label{ineq:dist:yap:yaq:1}
\lim_{\alpha\to +\infty}\frac{d_g(\yap,\yaq)}{\nap}=+\infty
\end{equation} 
\begin{equation}\label{ineq:dist:xai:yap:yaq:2}
\hbox{or }\left\{\lim_{\alpha\to +\infty}\frac{d_g(\yap,\yaq)}{\nap}=c_{p,q}\in (0,+\infty)\hbox{ and }\naq=o(\nap)\right\}.
\end{equation}
In addition, for any $p\in \{1,...,L\}$, there exists $V_p\in D_k^2(\rn)\cap C^{2k}(\rn)-\{0\}$ such that $\Delta_\xi^k V_p=|V_p|^{\crit-2}V_p$ strongly in $\rn$ and weakly in $D_k^2(\rn)$ and
\begin{equation}\label{cv:tvap}
\tvap :=\nap^{\frac{n-2k}{2}}\ua(\hbox{exp}_{\yap}(\nap \cdot))\to V_p\hbox{ in }C^{2k}_{loc}(\rn-S_p)
\end{equation} 
\begin{equation}\label{def:Sp}
\hbox{where }S_p:=\left\{\lim_{\alpha\to \infty}\frac{\hbox{exp}_{\yap}^{-1}(\xai)}{\nap}/\, i\in I_p\right\}.
\end{equation}
\end{defi}

\begin{proposition}\label{prop:4.4} Assume that $(\tilde{H}_L)$ holds for some $L\geq 0$. Then $K+L\leq \Lambda^{\crit} K(n,k)^{\frac{n}{2k}}$.
\end{proposition}

\begin{proof} It follows from \eqref{hyp:bnd:nrj:pf} that
\begin{equation}\label{eq:36}
\int_{\bigcup_i\Omega_{i,\alpha}(R)\cup \bigcup_p\tilde{\Omega}_{p,\alpha}(R) }|\ua|^{\crit}\, dv_g\leq \int_M |\ua|^{\crit}\, dv_g\leq\Lambda^{\crit}.
\end{equation}
It follows from \eqref{cv:tuai}, \eqref{cv:tvap} and \eqref{min:nrj:U} that for any $i\in\{1,...,K\}$ and $p\in\{1,...,L\}$, we have that
\begin{equation}
\left\{\begin{array}{c}
\lim_{R\to +\infty}\lim_{\alpha\to +\infty}\int_{\Omega_{i,\alpha}(R) }|\ua|^{\crit}\, dv_g\geq \frac{1}{K(n,k)^{\frac{n}{2k}}}\\
\lim_{R\to +\infty}\lim_{\alpha\to +\infty}\int_{\tilde{\Omega}_{p,\alpha}(R)}|\ua|^{\crit}\, dv_g\geq \frac{1}{K(n,k)^{\frac{n}{2k}}}\end{array}\right\}.\label{eq:37}
\end{equation}
Therefore, the conclusion of the proposition holds provided  that the integral is neglictible on the intersection of the domains. This is what we prove now.

\smallskip\noindent  For $i,j\in \{1,...,K\}$, $i\neq j$, it follows from \eqref{lim:d:infty} that for $\alpha>0$ large enough, we have that $\Omega_{i,\alpha}(R)\cap \Omega_{j,\alpha}(R)=\emptyset$. Therefore
\begin{equation}\label{lim:int:1}
\lim_{\alpha\to \infty}\int_{\Omega_{i,\alpha}(R)\cap \Omega_{j,\alpha}(R)}|\ua|^{\crit}\, dv_g=0.
\end{equation}
\smallskip\noindent For $i\in \{1,...,K\}$ and $p\in\{1,...,L\}$. Assume that, up to to extraction,
$$\Omega_{i,\alpha}(R)\cap \tilde{\Omega}_{p,\alpha}(R)\neq \emptyset\hbox{ for }\alpha\to +\infty.$$
Therefore, there exists $(a_\alpha)_\alpha\in M$ lying in the intersection, so that $d_g(a_\alpha, \xai)\leq R\mai$, $d_g(a_\alpha, \yap)\leq R\nap$ and $d_g(a_\alpha, \xaj)\geq R^{-1}\nap$ for all $j\in I_p$. The triangle inequality yields
\begin{equation}\label{eq:34}
d_g(\xai,\xaj)\geq d_g(a_\alpha,\xaj)-d_g(a_\alpha,\xai)\geq R^{-1}\nap-R\mai\hbox{ for all }j\in I_p.
\end{equation}
Another application of the triangle inequality yields $d_g(\xai,\yap)=O(\mai+\nap)$, and then, with  \eqref{ineq:dist:xai:yap}, we get 
\begin{equation}\label{eq:35}
\mai=o(\nap)\hbox{ and }d_g(\xai,\yap)=O(\nap).
\end{equation}
We set $\Theta_\alpha\in B_{\nap^{-1}i_g(M)}(0)\subset\rn$ be such that $\xai=\hbox{exp}_{\yap}(\nap \Theta_\alpha)$ for all $\alpha>0$. It follows from \eqref{eq:35} that $|\Theta_\alpha|=O(1)$ and then there exists $\Theta_\infty\in\rn$ such that $\lim_{\alpha\to +\infty}\Theta_\alpha=\Theta_\infty$. It follows from \eqref{lem:lip} that
$$\Omega_{i,\alpha}(R)\cap \tilde{\Omega}_{p,\alpha}(R)\subset B_{R\mai}(\xai)\subset \hbox{exp}_{\yap}\left(\nap B_{2R\mai/\nap}(\Theta_\alpha)\right).$$
We claim that $\Theta_\infty\not\in S_p$, where $S_p$ is as \eqref{def:Sp}. Indeed, it follows from  \eqref{eq:34} and \eqref{eq:35} that $d_g(\xai,\xaj)\geq (2R)^{-1}\nap$ for all $j\in I_p$, and then, using \eqref{lem:lip},  and passing to the limit, we get that $|\Theta_\infty-\Theta|\geq (4R)^{-1}$ for all $\Theta\in S_p$. Therefore $\Theta_\infty\not\in S_p$ and the claim is proved.

\smallskip\noindent Taking $\tvap$ as in \eqref{cv:tvap}, with a change of variable, we get that
\begin{eqnarray*}\int_{\Omega_{i,\alpha}(R)\cap \tilde{\Omega}_{p,\alpha}(R)}|\ua|^{\crit}\, dv_g&\leq& \int_{\hbox{exp}_{\yap}\left(\nap B_{2R\mai/\nap}\Theta_\alpha)\right)}|\ua|^{\crit}\, dv_g\\
&\leq& \int_{B_{2R\mai/\nap}(\Theta_\alpha)}|\tvap|^{\crit}\, dv_{\tilde{g}_{\alpha,p}}
\end{eqnarray*}
where $\tilde{g}_{\alpha,p}:=\hbox{exp}_{\yap}^{\star}g(\nap\cdot)$. It then follows from \eqref{cv:tvap}, $\Theta_\alpha\to\Theta_\infty\in\rn-S_p$ and  $\mai=o(\nap)$ that 
\begin{equation}\label{lim:int:2}
\lim_{\alpha\to \infty}\int_{\Omega_{i,\alpha}(R)\cap \tilde{\Omega}_{p,\alpha}(R)}|\ua|^{\crit}\, dv_g=0.
\end{equation}
Note that when the domain is empty, then this result is trivial.

\medskip\noindent We now choose $p,q\in \{1,...,L\}$ such that $p\neq q$. Arguing as above, we get that
\begin{equation}\label{lim:int:3}
\lim_{\alpha\to \infty}\int_{\tilde{\Omega}_{p,\alpha}(R)\cap \tilde{\Omega}_{q,\alpha}(R)}|\ua|^{\crit}\, dv_g=0.
\end{equation}


\smallskip\noindent We now can conclude the proof of Proposition \ref{prop:4.4}. It follows from \eqref{eq:36}, \eqref{eq:37}, \eqref{lim:int:1}, \eqref{lim:int:2} and \eqref{lim:int:3} that  $K+L\leq \Lambda^{\crit} K(n,k)^{\frac{n}{2k}}$. The proposition is proved.\end{proof}

\end{document}
